% =====================================================================
%  Part VIII -- ML framing and context.
%  Source: tier0.md Part VI (2852-3068) minus VI.6, which is relocated
%          to Part I; VI.7 items 1-5 (3096-3105), item 6 relocated to
%          Part I's self-check.
% =====================================================================
\part{ML framing and context}
\label{part:ml}

% ---------------------------------------------------------------------
\chapter{What kind of learning problem this is}
\label{sec:learningproblem}

\textbf{It is not neural operator learning in the DeepONet/FNO sense.} Those
target maps between \emph{function spaces} (infinite-dimensional ---
\citealt{lu2021}; \citealt{li2021}). Here input and output are fixed finite
vectors: \textbf{83 or 154 numbers in} --- ($\log T$, $\log\rho$, $\mathbf{X}$,
$\dd t$) per \S\ref{sec:onerecord} --- and \textbf{82 or 153 out} ---
($\mathbf{X}'$, $e_{\text{nuc}}$, $\varepsilon_\nu$). The ``operator'' framing
refers to the \emph{family} $\{\Phi_{\Delta t}\}$ indexed by timestep.

(Do not confuse this 83 with the $d \approx 82$ of \S\ref{sec:sobol}. That is
the \emph{Sobol sampling} dimension, 2 thermodynamic $+$ 80 composition, where
$\dd t$ is not a sampled coordinate but a nine-point grid applied to every drawn
state. Same-looking number, different object.)

\textbf{So what is hard about it?} Not expressivity --- a sufficiently large MLP
can represent the map. The difficulties are:

\begin{enumerate}
\item \textbf{Dynamic range.} $X$ spans 15 decades; $\varphi$ spans more.
      Standard initialization and standard losses both implicitly assume $O(1)$
      targets.
\item \textbf{Stiffness in the target.} Neighbouring inputs can map to very
      different outputs when a state is near a fast-mode threshold --- the
      learned map must have large, spatially varying Lipschitz constants, which
      is exactly what smooth function approximators resist.
\item \textbf{Exact constraints.} No amount of capacity produces exact
      conservation; it must be architectural
      (\S\ref{sec:conservationfailure}, \S\ref{sec:whereconservation}).
\item \textbf{Extrapolation under rollout.} The deployment distribution is the
      reachable manifold (\S\ref{sec:onezone}), and rollout drifts off it.
\end{enumerate}

Recognizing that expressivity is \emph{not} the bottleneck is what justifies
spending the effort on the decoder structure rather than on backbone scale.

% ---------------------------------------------------------------------
\chapter{Message passing, concretely}
\label{sec:messagepassing}

A GNN layer updates node states by aggregating over neighbours (the
message-passing scheme of \citealt{gilmer2017}; \citealt{battaglia2018}):

\begin{equation}
h_i^{(k+1)} = \phi\Bigl(h_i^{(k)},\; \bigoplus_{j \in N(i)}
  \psi\bigl(h_i^{(k)}, h_j^{(k)}, e_{ij}\bigr)\Bigr)
\label{eq:message-passing}
\end{equation}

with $\bigoplus$ a permutation-invariant aggregator (sum, mean, max).

\textbf{Bipartite version}, which is what this network is:

\begin{align*}
\text{reaction node:}\quad h_r &\leftarrow
  \psi_{IR}\bigl(\{h_i : i \in \text{reactants}(r)\},\; \text{edge features}\bigr)\\
\text{species node:}\quad h_i &\leftarrow
  \psi_{RI}\Bigl(h_i,\; \bigoplus_{r \ni i} (\nu_{ir},\, h_r)\Bigr)
\end{align*}

One full round $=$ species $\to$ reaction $\to$ species. \textbf{Sum aggregation
is the physically correct choice on the R$\to$I half-step}, because the true
update is literally a signed sum: $\dd Y_i = \sum_r \nu_{ir} \varphi_r$. Mean or
max would destroy extensivity (and sum is independently the most expressive
aggregator in the multiset sense --- \citealt{xu2019}).

\textbf{Depth --- and state the hop convention, or the arithmetic does not
close.} One full round is species $\to$ reaction $\to$ species $=$ \textbf{2
bipartite hops}, so $K$ rounds propagate information $2K$ bipartite hops, not
$K$. With measured bipartite radius 3 and diameter 6:

\begin{equation*}
2K \geq \text{diameter} = 6 \quad\Longrightarrow\quad K \geq 3
\end{equation*}

So \textbf{$K = 3$ rounds suffice to span the graph}, and the project's $K = 5$
carries ${\sim}4$ hops of margin. (The derived $I\to I$ projection does
\emph{not} simply halve these numbers: the repo measures it at radius 2,
diameter 4 --- \code{RESULTS.md} 2026-07-09 --- because the projection is its
own graph object, not the bipartite metric divided by two. At one $I\to I$ hop
per round, that view requires $K \geq 4$. Either way, $K = 5$ exceeds the
requirement.) Beware the version of this claim that reads ``radius 3, diameter 6
$\implies K = 5$'': under a one-hop-per-round convention $K = 5$ would
\emph{fail} to cover a diameter-6 graph, and under the two-hop convention it
overshoots. Neither reading yields 5 --- the number is a design margin, and
should be defended as one.

Deeper is not obviously better:

\begin{itemize}
\item \textbf{Oversmoothing:} as $K$ grows, node representations converge toward
      a common value and become indistinguishable \citep{li2018}. At $K = 5$ not
      yet serious, but it bounds how much depth can help.
\item \textbf{Expressivity:} message passing is bounded above by the
      1-Weisfeiler--Lehman test in distinguishing power
      (\citealt{xu2019}; \citealt{morris2019} --- the theorems hold for labelled
      graphs too; 1-WL runs on the initial node labels). Irrelevant here ---
      nodes carry distinguishing features ($Z$, $A$, $Y$), so the colouring
      separates nodes from iteration 0 and the bound never binds.
\end{itemize}

\textbf{Equivariance.} The physics is invariant to relabelling species.
Equation \eqref{eq:message-passing} is permutation-equivariant by construction
\citep{battaglia2018,bronstein2021}. A dense MLP must learn this from data,
spending capacity on a symmetry that could have been free.

% ---------------------------------------------------------------------
\chapter{Feature and loss design under 15 decades of dynamic range}
\label{sec:featureloss}

\section{Inputs first --- the transform that is legal}
\label{sec:inputtransform}

The raw composition vector spans 15 decades, which no standard initialization or
normalization scheme tolerates. The conventional fix is a signed-log or
inverse-hyperbolic-sine transform on the \emph{inputs}:

\begin{equation*}
\tilde{x}_i = \asinh\!\left(\frac{X_i}{s}\right) \qquad \text{or} \qquad
\log_{10}(X_i + X_0)
\end{equation*}

\textbf{This is legal, and it is important to be clear about why}, given
\S\ref{sec:whereconservation} spends its length forbidding exactly this
function. The rule is not ``asinh is banned''. The rule is:

\begin{invariant}
asinh/signed-log are fine as \textbf{internal latents}, including on inputs.
They may never be the space in which a \textbf{sum constraint or projection} is
evaluated.
\end{invariant}

Encoding $X$ into a latent and decoding $\varphi$ (or $\dd Y$) in a linear space
keeps the conservation map downstream of every nonlinearity, which is the whole
requirement. The NuGNN failure mode was subtler than ``projecting in the warped
space'': because its outputs live in a signed-log space, any explicit
conservation enforcement (loss term or output activation) required inverting the
warp inside training, which destabilized optimization --- so enforcement was
dropped entirely, leaving an unconstrained model plus a retry-on-drift heuristic
at rollout \citep{kim2026}. The lesson is the same, not the warp itself: a sum
constraint is only enforceable where the output is linear.

Other input channels worth having, most of them free:

\begin{itemize}
\item ($\log T$, $\log\rho$) --- already logarithmic, already $O(1)$ after box
      normalization
\item per-node static features: $Z$, $A$, $N-Z$, binding energy per nucleon, the
      weak/strong flag
\item per-node dynamic features: $Y_i$, and the equilibrium departure $\delta_i$
      if computed
\item per-reaction features: $Q$-value, chapter/$\Delta N$, screening factor,
      and \textbf{$\kappa$ diagnostics}
\end{itemize}

That last one is a project decision worth understanding: once the Guidry mask
was measured empty, the $\kappa$ information did not become useless --- it was
demoted from a \emph{mask} (a hard architectural gate) to a \emph{feature}
(something the model may condition on). Carrying a diagnostic as a feature is
the fallback when it is informative but not cleanly separable.

\textbf{The zero problem.} Many $X_i$ are exactly 0 in initial conditions, and
$\log(0)$ is undefined. asinh handles it ($\asinh 0 = 0$); log needs an additive
floor, which then becomes a hyperparameter encoding what counts as negligible.
Prefer asinh on inputs for this reason alone.

\section{Then losses --- where the naive choices all fail}
\label{sec:losses}

\begin{center}
\small
\begin{tabularx}{\textwidth}{@{}R{0.75} R{1.25}@{}}
\toprule
\textbf{Loss} & \textbf{Failure mode} \\
\midrule
MSE on $X$ &
  Dominated entirely by the few species with $X \sim 1$; trace species
  contribute nothing \\
MSE on $\log X$ &
  Undefined at $X = 0$; and the $10^{-15}$ floor is \emph{censored}, so
  log-space error there measures a clamp \\
Relative error $|\Delta X|/X$ &
  Explodes for trace species; amplifies label noise where $X$ is at the float32
  noise level \\
Relative with floor $|\Delta X|/(X + X_0)$ &
  Workable --- but $X_0$ is now a hyperparameter encoding what you consider
  negligible \\
\bottomrule
\end{tabularx}
\end{center}

\textbf{The structural observations that reframe the problem:}

\begin{itemize}
\item The quantity that must be accurate is $\Ye = \sum Z_i Y_i$, a weighted
      sum. Errors in species with small $Z_i/A_i$ deviation from 0.5 barely
      matter. A loss weighted toward \Ye{}-controlling species is physically
      motivated, not a hack --- and the project already identifies the target
      channels (top-$|\dd\dot{Y}_e|$ weak channels: \nuc{Ni}{56} EC,
      \nuc{S}{31} EC, \nuc{Fe}{52} EC, $p$ EC).
\item Supervising in $\Phi$ \textbf{space} rather than $\Delta X$ space
      side-steps the dynamic-range problem partly, because fluxes are more
      nearly log-uniform. This is exactly what the auxiliary $\Phi$ labels from
      the reference integrator are for.
\item \textbf{Censoring must be handled explicitly.} A value at $10^{-15}$ means
      ``$\leq 10^{-15}$'', not ``$= 10^{-15}$''. The correct treatment is a
      one-sided (censored) likelihood --- the standard Tobit treatment
      \citep{tobin1958} --- not a squared error against the clamp.
\end{itemize}

% ---------------------------------------------------------------------
\chapter{MESA and bbq, concretely}
\label{sec:mesabbq}

\textbf{MESA} (Modules for Experiments in Stellar Astrophysics) is a modular 1-D
stellar evolution code \citep{paxton2011,jermyn2023}. The relevant modules:

\begin{itemize}
\item \code{net} --- the nuclear reaction network solver
\item \code{rates} --- REACLIB fits, tabulated weak rates, screening
\item \code{weaklib} --- tabulated weak rates on $(T, \rho\Ye)$ grids, with
      source precedence LMP $>$ \citeauthor{oda1994} $>$ \citeauthor{ffn} (\citealt{paxton2011}; in r23.05.1
      this machinery lives inside the \code{rates} module rather than as a
      separate directory)
\end{itemize}

\textbf{Network types} in MESA (``hard-wired'' is MESA's own terminology;
``softwired'' is community usage --- \citealt{farmer2016} classify the
\code{mesa\_*.net} family exactly this way):

\begin{center}
\small
\begin{tabularx}{\textwidth}{@{}l Y@{}}
\toprule
\textbf{Type} & \textbf{Meaning} \\
\midrule
\emph{hardwired} & Fixed species and fixed reaction list, compiled in \\
\textbf{softwired} &
  Fixed species list; reaction links assembled at runtime from whatever rates
  connect them --- \textbf{what \msn{80}/\msn{151} are} \\
\emph{approx} &
  Reduced networks with lumped/approximated links (e.g.\ \code{approx21}) \\
\bottomrule
\end{tabularx}
\end{center}

A \code{.net} file declares the species; the links follow. This is why the
reaction \emph{count} (607 / 1518) is something you have to \emph{measure} by
asking MESA what it built, rather than read off a list --- the entire point of
the Fortran probe in \code{src/mesa\_probes/probe.f90} and the reconciliation in
S6.

\textbf{bbq} is a one-zone burner utility built on MESA's \code{net}
(\citealt{farmerbbq}, Zenodo record 7585202; the tool behind
\citealt{grichener2025}): give it $(T, \rho, \mathbf{X}_0)$ and a time grid, get
back the composition history and energetics. Its \code{use\_hydrostatic} $+$
\code{times\_from\_file} mode produced both the shipped test trajectories and
the local rerun campaign.

% ---------------------------------------------------------------------
\chapter{Prior art in network acceleration --- where ML actually sits}
\label{sec:priorart}

ML is the newest entry in a long line of attempts. Knowing the line matters,
because the project's vocabulary is borrowed from it:

\begin{center}
\small
\begin{tabularx}{\textwidth}{@{}R{0.8} R{1.25} R{0.95}@{}}
\toprule
\textbf{Approach} & \textbf{Idea} &
\textbf{Why it appears in this project} \\
\midrule
\textbf{NSE tabulation} &
  Above ${\sim}7$~GK, skip the ODE --- composition is algebraic in
  $(T,\rho,\Ye)$ \citep{hix1999b} &
  Sets the box's upper edge (\S\ref{sec:regimebox}) \\
\textbf{QSE / cluster reduction} \citep{hix1996,hix1999a,hix2007} &
  Equilibrated groups collapse to one unknown per cluster; integrate only
  bridges $+$ weak &
  The \code{qse/} solver, the Si-group config, \code{r\_QSE} \\
\textbf{Partial equilibrium / explicit asymptotic} \citep{guidry2013} &
  Detect equilibrated forward--reverse \emph{pairs} dynamically, drop them from
  integration, use explicit methods for the rest &
  The ``Guidry mask'', the $\delta_r$ criterion, the $\epsilon$ sweep \\
\textbf{Adaptive networks} &
  Grow/shrink the species set by local importance &
  Related to the active-set framing of the kill-test \\
\textbf{ML surrogates} & Learn the flow map &
  This project, and the NNN baseline \\
\bottomrule
\end{tabularx}
\end{center}

\begin{aside}
\textbf{Read the kill-test as a test of whether the QSE-reduction idea applies
to \emph{this} data.} The Guidry mask being measured \emph{empty} is a statement
that the classical reduction strategy does not have a clean target here ---
itself a result, and why the flux head ended up full-width.
\end{aside}

% ---------------------------------------------------------------------
\chapter*{Self-check --- ML framing}
\addcontentsline{toc}{chapter}{Self-check --- ML framing}
\label{sec:selfcheck-ml}

\begin{selfcheck}
\item Explain why this is not neural-operator learning in the DeepONet/FNO
      sense, and what the ``operator'' framing does refer to.
\item Write \eqref{eq:message-passing} for the bipartite network and justify sum
      aggregation on the R$\to$I half-step from equation
      \eqref{eq:network-ode}.
\item Give a concrete failure case for each of the four losses in
      \S\ref{sec:losses}, using real numbers from the box ($X \sim 1$ species vs
      $X \sim 10^{-12}$ species).
\item Explain why the reaction count must be \emph{measured} in a softwired
      network, and name the code path that measures it.
\item For each entry in the \S\ref{sec:priorart} prior-art table, name the module
      or config in this repository that inherits its vocabulary.
\end{selfcheck}

\begin{aside}
\textbf{One question from the source document's \S VI.7 is not here.} Item~6
(tracing the inference chain from a 1\% \Ye{} error to a changed observable)
tests \S VI.6, whose material now opens the report as
\S\ref{sec:obsgrounding}; it appears in Part~\ref{part:context}'s self-check
instead. See Appendix~\ref{app:concordance}.
\end{aside}
